\documentclass[aspectratio=169,11pt]{beamer}
\usetheme{default}
\usefonttheme{professionalfonts}

\usepackage[utf8]{inputenc}
\usepackage[spanish]{babel}
\usepackage[T1]{fontenc}
\usepackage{graphicx}
\usepackage{hyperref}
\usepackage{booktabs}
\usepackage{array}
\usepackage{tabularx}
\usepackage{xcolor}
\usepackage{tikz}
\usetikzlibrary{positioning,arrows.meta,shapes.geometric,fit,calc}

% Colores
\definecolor{paperbg}{RGB}{250,250,248}
\definecolor{darkred}{RGB}{165,36,52}
\definecolor{accentteal}{RGB}{21,112,112}
\definecolor{accentgold}{RGB}{183,125,36}
\definecolor{accentblue}{RGB}{49,100,168}
\definecolor{ink}{RGB}{28,35,43}
\definecolor{muted}{RGB}{88,99,111}
\definecolor{softred}{RGB}{250,231,234}
\definecolor{softblue}{RGB}{232,239,250}
\definecolor{softgreen}{RGB}{224,242,238}
\definecolor{softgold}{RGB}{251,241,220}
\definecolor{softgray}{RGB}{242,244,246}
\definecolor{linegray}{RGB}{214,221,228}

\setbeamercolor{normal text}{fg=ink,bg=paperbg}
\setbeamercolor{title}{fg=ink}
\setbeamercolor{subtitle}{fg=muted}
\setbeamercolor{frametitle}{fg=ink,bg=paperbg}
\setbeamercolor{block title}{fg=ink,bg=paperbg}
\setbeamercolor{block body}{bg=white,fg=ink}
\setbeamerfont{title}{size=\Huge,series=\bfseries}
\setbeamerfont{subtitle}{size=\large}
\setbeamerfont{frametitle}{size=\Large,series=\bfseries}
\setbeamerfont{block title}{size=\normalsize,series=\bfseries}
\setbeamerfont{block body}{size=\normalsize}
\setbeamersize{text margin left=8mm,text margin right=8mm}
\setbeamertemplate{navigation symbols}{}
\setbeamertemplate{frametitle}{%
  \vspace{0.8mm}
  \begin{tikzpicture}[remember picture,overlay]
    \fill[darkred] (current page.north west) rectangle ([xshift=2.0mm,yshift=-\paperheight]current page.north west);
    \fill[accentteal] ([xshift=2.0mm]current page.north west) rectangle ([xshift=3.2mm,yshift=-\paperheight]current page.north west);
  \end{tikzpicture}
  \insertframetitle\par
  \vspace{0.2mm}
  {\color{linegray}\rule{\linewidth}{0.45pt}}
}
\setbeamertemplate{block begin}{%
  \par\vskip0.35ex
  \begin{beamercolorbox}[wd=\linewidth,sep=0.58em,rounded=true,shadow=false]{block body}
    {\usebeamerfont{block title}\color{darkred}\insertblocktitle}\par\vspace{0.25em}
    \usebeamerfont{block body}
}
\setbeamertemplate{block end}{%
  \end{beamercolorbox}
  \vskip0.45ex
}
\setbeamertemplate{footline}{%
  \pgfmathsetlengthmacro{\progressbar}{\paperwidth*\insertframenumber/\inserttotalframenumber}
  \begin{tikzpicture}[remember picture,overlay]
    \fill[linegray] ([yshift=2.3mm]current page.south west) rectangle ([yshift=2.55mm]current page.south east);
    \fill[darkred] ([yshift=2.3mm]current page.south west) rectangle
      ([xshift=\progressbar,yshift=2.55mm]current page.south west);
  \end{tikzpicture}
  \hfill{\scriptsize\color{muted}\insertframenumber/\inserttotalframenumber}\hspace*{8mm}\vspace{2.0mm}
}

\newcommand{\tagbox}[2]{%
  \tikz[baseline=(X.base)]\node[rounded corners=9pt,inner xsep=6pt,inner ysep=3pt,fill=#1,draw=linegray] (X) {\scriptsize #2};%
}
\newcommand{\kicker}[1]{{\small\bfseries\color{accentteal}#1}}
\newcommand{\lead}[1]{{\large\color{ink}#1}}
\newcommand{\smallmuted}[1]{{\scriptsize\color{muted}#1}}
\newcommand{\metriccard}[3]{%
  \begin{beamercolorbox}[wd=\linewidth,sep=0.58em,rounded=true,shadow=false]{block body}
    {\scriptsize\color{muted}#1}\par
    {\Large\bfseries\color{#2}#3}
  \end{beamercolorbox}
}

\title{Moderador de Contenido en Videos de YouTube}
\subtitle{Corpus peruano, etiquetado humano y clasificador local de riesgo textual}
\author{
  \textbf{Grupo 4}\\
  Koc Góngora, Luis Enrique\\
  Mancilla Antay, Alex Felipe\\
  Melendez Garcia, Herbert Antonio\\
  Paitan Cano, Dennis Jack
}
\institute{
Maestría en Inteligencia Artificial\\
Universidad Nacional de Ingeniería (UNI)\\
Procesamiento de Lenguaje Natural
}
\date{Semestre 2026-1}

\begin{document}

% ==================== 1 ====================
\begin{frame}[plain]
  \begin{tikzpicture}[remember picture,overlay]
    \fill[paperbg] (current page.south west) rectangle (current page.north east);
    \fill[darkred] (current page.north west) rectangle ([xshift=3.4mm]current page.south west);
    \fill[accentteal] ([xshift=3.4mm]current page.north west) rectangle ([xshift=5.1mm]current page.south west);
    \fill[softred] ([xshift=9.5cm,yshift=-0.55cm]current page.north west) rectangle ([xshift=15.9cm,yshift=-1.15cm]current page.north west);
    \fill[softgreen] ([xshift=9.5cm,yshift=-1.35cm]current page.north west) rectangle ([xshift=14.8cm,yshift=-1.95cm]current page.north west);
    \fill[softblue] ([xshift=9.5cm,yshift=-2.15cm]current page.north west) rectangle ([xshift=15.4cm,yshift=-2.75cm]current page.north west);
  \end{tikzpicture}

  \vspace{8mm}
  \kicker{Procesamiento de Lenguaje Natural}

  \vspace{7mm}
  {\usebeamerfont{title}\color{ink}\inserttitle\par}
  \vspace{2mm}
  {\usebeamerfont{subtitle}\color{muted}\insertsubtitle\par}

  \vspace{7mm}
  \tagbox{softgreen}{corpus peruano}
  \tagbox{softblue}{etiquetado humano}
  \tagbox{softred}{clasificador local}

  \vfill
  {\small\color{ink}\textbf{Grupo 4}}\par
  {\scriptsize\color{muted}Koc Góngora, Luis Enrique \quad Mancilla Antay, Alex Felipe\par
  Melendez Garcia, Herbert Antonio \quad Paitan Cano, Dennis Jack}\par
  \vspace{2mm}
  {\scriptsize\color{muted}Maestría en Inteligencia Artificial -- UNI \quad | \quad Semestre 2026-1}
\end{frame}

% ==================== 2 ====================
\begin{frame}
  \frametitle{Idea Central del Trabajo}

  \lead{El trabajo construye un moderador textual para videos peruanos, con corpus propio, anotación humana y evidencia visible para revisión.}

  \vspace{5mm}
  \begin{columns}[T]
    \column{0.31\textwidth}
      \metriccard{Objeto de estudio}{accentblue}{Videos públicos}
      \vspace{1mm}
      {\small Política y farándula analizadas desde título, descripción, subtítulos y transcripción.}

    \column{0.31\textwidth}
      \metriccard{Problema aplicado}{darkred}{Revisión lenta}
      \vspace{1mm}
      {\small El volumen crece y el lenguaje depende del contexto peruano, con ironía y modismos.}

    \column{0.31\textwidth}
      \metriccard{Artefacto}{accentteal}{Flujo local}
      \vspace{1mm}
      {\small Scraping, chunks, etiquetado, entrenamiento e inferencia sin servicios externos.}
  \end{columns}

  \vfill
  \centering
  \tagbox{softgreen}{sin APIs externas}
  \tagbox{softblue}{datos trazables}
  \tagbox{softred}{decisión con revisión humana}
\end{frame}

% ==================== 3 ====================
\begin{frame}
  \frametitle{Contexto y Problemática}

  \kicker{Dónde aparece la dificultad}
  \vspace{1.5mm}

  \begin{columns}[T]
    \column{0.32\textwidth}
      \metriccard{Fuente}{accentblue}{YouTube}
      \vspace{1mm}
      {\small Opinión política, noticias, entrevistas, clips y entretenimiento en un flujo continuo.}

    \column{0.32\textwidth}
      \metriccard{Lenguaje}{accentgold}{Local}
      \vspace{1mm}
      {\small Modismos, ironía, apodos y referencias culturales que cambian el sentido del texto.}

    \column{0.32\textwidth}
      \metriccard{Riesgo}{darkred}{No auditable}
      \vspace{1mm}
      {\small Un dictamen automático sin fragmentos activadores no permite revisar ni corregir.}
  \end{columns}

  \vspace{7mm}
  \centering
  \begin{tikzpicture}[
    box/.style={draw=linegray,rounded corners=10pt,fill=white,minimum width=2.85cm,minimum height=0.78cm,align=center,font=\small},
    arr/.style={-Latex,thick,draw=darkred}
  ]
    \node[box,fill=softblue] (a) {Videos públicos};
    \node[box,right=0.45cm of a,fill=softgray] (b) {Texto transcrito};
    \node[box,right=0.45cm of b,fill=softred] (c) {Riesgo textual};
    \node[box,right=0.45cm of c,fill=softgreen] (d) {Revisión humana};
    \draw[arr] (a) -- (b);
    \draw[arr] (b) -- (c);
    \draw[arr] (c) -- (d);
  \end{tikzpicture}
\end{frame}

% ==================== 4 ====================
\begin{frame}
  \frametitle{Brecha y Formulación del Problema}

  \begin{block}{Brecha}
    Existen clasificadores de toxicidad y odio, pero el proyecto requiere un corpus local, una guía de anotación sencilla y un flujo reproducible que funcione sin servicios externos.
  \end{block}

  \begin{columns}[T]
    \column{0.52\textwidth}
      \begin{block}{Problema general}
        ¿Cómo construir un moderador textual local que clasifique fragmentos de videos peruanos y entregue evidencia verificable para apoyar la revisión humana?
      \end{block}

    \column{0.42\textwidth}
      \begin{block}{Problemas específicos}
        \begin{itemize}
          \item ¿Qué fuentes y categorías usar?
          \item ¿Cómo crear chunks auditables?
          \item ¿Cómo entrenar un baseline defendible?
          \item ¿Cómo agregar riesgo por video?
        \end{itemize}
      \end{block}
  \end{columns}

\end{frame}

% ==================== 5 ====================
\begin{frame}
  \frametitle{Propuesta de Solución: Artefacto}

  \begin{center}
  \begin{tikzpicture}[
    stage/.style={draw=linegray,rounded corners,fill=white,minimum width=1.95cm,minimum height=0.9cm,align=center,font=\small},
    artifact/.style={draw=darkred,thick,rounded corners,fill=softred,minimum width=2.25cm,minimum height=0.95cm,align=center,font=\small\bfseries},
    arr/.style={-Latex,thick,draw=darkred}
  ]
    \node[stage,fill=softblue] (s1) {Scraping\\web público};
    \node[stage,right=0.22cm of s1] (s2) {Subtítulos o\\ASR local};
    \node[stage,right=0.22cm of s2] (s3) {Limpieza y\\chunks};
    \node[stage,right=0.22cm of s3] (s4) {Etiquetado\\HTML};
    \node[stage,right=0.22cm of s4] (s5) {Modelo\\local};
    \node[artifact,right=0.22cm of s5] (s6) {Frontend\\producción};
    \draw[arr] (s1) -- (s2);
    \draw[arr] (s2) -- (s3);
    \draw[arr] (s3) -- (s4);
    \draw[arr] (s4) -- (s5);
    \draw[arr] (s5) -- (s6);
  \end{tikzpicture}
  \end{center}

  \vspace{3mm}
  \begin{block}{Producto verificable}
    El artefacto no reemplaza al moderador humano. Prioriza fragmentos, asigna categorías sugeridas, muestra texto activador y permite auditar errores.
  \end{block}

  \begin{block}{Restricción de implementación}
    La recolección, etiquetado, entrenamiento e inferencia se ejecutan localmente. No se usan APIs para obtener etiquetas ni para clasificar contenido.
  \end{block}
\end{frame}

% ==================== 6 ====================
\begin{frame}
  \frametitle{Canales Sugeridos para Scraping}

  \lead{La recolección debe mezclar fuentes de lenguaje formal, opinión, humor, farándula y viajes para que el clasificador no aprenda un solo registro.}

  \vspace{5mm}
  \begin{columns}[T]
    \column{0.24\textwidth}
      \metriccard{01}{accentblue}{Formal}
      \vspace{1mm}
      {\scriptsize Noticias, entrevistas y periodismo con dicción más estable.}

    \column{0.24\textwidth}
      \metriccard{02}{accentteal}{Opinión}
      \vspace{1mm}
      {\scriptsize Debate político, sátira y explicaciones con lenguaje local.}

    \column{0.24\textwidth}
      \metriccard{03}{darkred}{Farándula}
      \vspace{1mm}
      {\scriptsize Conflicto público, rumores, ataques personales y contexto sensible.}

    \column{0.24\textwidth}
      \metriccard{04}{accentgold}{Viajes}
      \vspace{1mm}
      {\scriptsize Lenguaje descriptivo y limpio para balancear clases seguras.}
  \end{columns}

  \vfill
  \tagbox{softgreen}{priorizar subtítulos}
  \tagbox{softblue}{audio claro}
  \tagbox{softgold}{episodios completos}
  \tagbox{softred}{validar canal oficial}
\end{frame}

% ==================== 7 ====================
\begin{frame}
  \frametitle{Candidatos Principales para Descarga}

  \kicker{Lista operativa alineada con el cuaderno 01}
  \vspace{2mm}

  \begin{columns}[T]
    \column{0.33\textwidth}
      \metriccard{Política}{accentblue}{Análisis}
      \vspace{1mm}
      {\small
      \begin{itemize}
        \item Marco Sifuentes / Ocram
        \item El diario de Curwen
        \item Sin Guion / Rosa María Palacios
        \item RPP, Canal N, Exitosa
        \item Willax y opinión
      \end{itemize}}

    \column{0.33\textwidth}
      \metriccard{Humor / streaming}{accentgold}{Contexto}
      \vspace{1mm}
      {\small
      \begin{itemize}
        \item Hablando Huevadas
        \item Todo Good
        \item Clips de entrevistas
        \item Conversación informal
      \end{itemize}}

    \column{0.33\textwidth}
      \metriccard{Farándula}{darkred}{Contraste}
      \vspace{1mm}
      {\small
      \begin{itemize}
        \item Magaly TV La Firme
        \item Amor y Fuego
        \item América Hoy
        \item Programas de espectáculo
      \end{itemize}}
  \end{columns}
\end{frame}

% ==================== 8 ====================
\begin{frame}
  \frametitle{Candidatos Complementarios y Verificación}

  \begin{columns}[T]
    \column{0.33\textwidth}
      \metriccard{Por confirmar}{darkred}{URL oficial}
      \vspace{1mm}
      {\small
      \begin{itemize}
        \item Cacas / El Cacas
        \item Goblinciano
        \item Instarándula
        \item Otros clips virales
      \end{itemize}}

    \column{0.33\textwidth}
      \metriccard{Viajes}{accentteal}{Clase segura}
      \vspace{1mm}
      {\small
      \begin{itemize}
        \item Misias pero viajeras
        \item Buen Viaje
        \item Viaja y Prueba
        \item Canales regionales
      \end{itemize}}

    \column{0.33\textwidth}
      \begin{block}{Regla de descarga}
        Usar solo canales públicos con URL oficial verificada, audio claro y episodios completos. Registrar descartes por copyright, duplicados o baja calidad.
      \end{block}
  \end{columns}
\end{frame}

% ==================== 8A ====================
\begin{frame}
  \frametitle{Taxonomía desde la Literatura Peruana}

  \lead{Las cuatro categorías base requieren subcategorías porque el racismo, el acoso y el daño en Perú tienen formas específicas que taxonomías genéricas no capturan.}

  \vspace{4mm}
  \begin{columns}[T]
    \column{0.32\textwidth}
      \begin{block}{Racismo negado}
        {\small Portocarrero (2009): ``fundamento invisible''. El racismo se practica pero se niega; aparece como criterio de educación o clase.}
      \end{block}

    \column{0.32\textwidth}
      \begin{block}{Racismo lingüístico}
        {\small Almeida \& Zavala (2022): burla del acento andino, el ``motoseo'' y la ortografía de migrantes como formas de discriminación en redes.}
      \end{block}

    \column{0.32\textwidth}
      \begin{block}{Falla en idiomas locales}
        {\small Thakur / CDT (2025): los sistemas automáticos fallan con español andino, quechua mezclado y modismos regionales.}
      \end{block}
  \end{columns}

  \vfill
  \tagbox{softred}{no transferible del inglés}
  \tagbox{softblue}{humor como cobertura del daño}
  \tagbox{softgreen}{etiquetado humano esencial}
\end{frame}

% ==================== 8B ====================
\begin{frame}
  \frametitle{Categorías y Subcategorías de Etiquetado}

  \kicker{Categorías no mutuamente excluyentes · 12 subcategorías con fundamento peruano}
  \vspace{2mm}

  \begin{columns}[T]
    \column{0.24\textwidth}
      \metriccard{Clase 0}{accentteal}{Seguro}
      \vspace{1mm}
      {\scriptsize
        \textbf{seguro}\\sin ataque directo\\
        \vspace{1mm}
        \textbf{seguro\_ironía\_marcada}\\parodia señalizada
      }

    \column{0.24\textwidth}
      \metriccard{Clase 1}{accentblue}{Racismo}
      \vspace{1mm}
      {\scriptsize
        étnico\_explícito\\linguístico\\
        clasismo\_racial\\regional\\encubierto\\
        {\tiny (Zavala \& Back, 2017;\\Callirgos, 1993)}
      }

    \column{0.24\textwidth}
      \metriccard{Clase 2}{darkred}{Acoso}
      \vspace{1mm}
      {\scriptsize
        misoginia\_género\\homofobia\\
        personal\_doxeo\\amenaza\_directa\\
        {\tiny (Monge-Olivarría \&\\Guerra-Corrales, 2023)}
      }

    \column{0.24\textwidth}
      \metriccard{Clase 3}{accentgold}{Sexual}
      \vspace{1mm}
      {\scriptsize
        explícito\\cosificación\\
        no\_consensual
      }
  \end{columns}

  \vfill
  \tagbox{softgreen}{12 subcategorías}
  \tagbox{softblue}{multi-etiqueta permitido}
  \tagbox{softred}{basadas en expertos peruanos}
\end{frame}

% ==================== 9 ====================
\begin{frame}
  \frametitle{Flags Transversales: Ironía y Humor como Cobertura}

  \lead{No son categorías de daño sino indicadores de ambigüedad que el moderador humano debe resolver antes de confirmar la etiqueta.}

  \vspace{4mm}
  \begin{columns}[T]
    \column{0.32\textwidth}
      \begin{block}{ironia\_ambigua}
        {\small No se determina si el daño es intencional o parodia crítica del racismo. Frecuente en canales de opinión (Curwen, Goblinciano).\\{\scriptsize (Vich, 2018)}}
      \end{block}

    \column{0.32\textwidth}
      \begin{block}{humor\_encubridor}
        {\small El hablante usa el humor para negar el daño (``es broma''). Brañez Medina (2012): ``racismo negado en la vida cotidiana''. El daño puede ser real.}
      \end{block}

    \column{0.32\textwidth}
      \begin{block}{contexto\_necesario}
        {\small El chunk aislado no permite clasificar; se requiere el video completo. Ocurre con español andino y modismos regionales.\\{\scriptsize (Thakur / CDT, 2025)}}
      \end{block}
  \end{columns}

  \vfill
  \tagbox{softgold}{coexisten con una categoría de daño}
  \tagbox{softblue}{activan revisión humana obligatoria}
  \tagbox{softgreen}{guía: 7 pasos, 13 ejemplos anotados}
\end{frame}

\begin{frame}
  \frametitle{Control de Calidad Humana}

  \begin{columns}[T]
    \column{0.31\textwidth}
      \metriccard{Muestra}{accentblue}{2 anotadores}
      \vspace{1mm}
      {\small Revisión independiente de una parte del corpus.}

    \column{0.31\textwidth}
      \metriccard{Consenso}{accentgold}{Casos dudosos}
      \vspace{1mm}
      {\small Discusión documentada para ejemplos ambiguos.}

    \column{0.31\textwidth}
      \metriccard{Acuerdo}{accentteal}{Kappa}
      \vspace{1mm}
      {\small Reporte de Cohen's kappa antes del entrenamiento final.}
  \end{columns}

  \vspace{5mm}
  \begin{block}{Regla metodológica}
    Antes de ampliar etiquetas, se revisan errores, desacuerdos y sesgos por tipo de canal. La guía se actualiza con ejemplos reales.
  \end{block}
\end{frame}

% ==================== 8 ====================
\begin{frame}
  \frametitle{Unidad de Análisis}

  \begin{center}
  \begin{tikzpicture}[
    box/.style={draw=linegray,rounded corners,fill=white,minimum width=2.9cm,minimum height=0.8cm,align=center,font=\small},
    chunk/.style={draw=linegray,rounded corners,fill=softblue,minimum width=1.8cm,minimum height=0.65cm,align=center,font=\scriptsize},
    arr/.style={-Latex,draw=darkred,thick}
  ]
    \node[box,fill=softgray] (video) {Video público};
    \node[chunk,below left=0.95cm and -0.15cm of video] (c1) {chunk 1};
    \node[chunk,right=0.28cm of c1] (c2) {chunk 2};
    \node[chunk,right=0.28cm of c2] (c3) {chunk 3};
    \node[chunk,right=0.28cm of c3] (c4) {chunk n};
    \node[box,below=1.0cm of c2,fill=softred] (clf) {Clasificador};
    \node[box,right=1.05cm of clf,fill=softgreen] (out) {Riesgo por video};
    \draw[arr] (video) -- (c2);
    \draw[arr] (c1) -- (clf);
    \draw[arr] (c2) -- (clf);
    \draw[arr] (c3) -- (clf);
    \draw[arr] (c4) -- (clf);
    \draw[arr] (clf) -- (out);
  \end{tikzpicture}
  \end{center}

  \vspace{1mm}
  \begin{block}{Qué se etiqueta}
    La unidad mínima es el fragmento textual. El video conserva metadatos, fuente y evidencias para auditar el resultado agregado.
  \end{block}
\end{frame}

\begin{frame}
  \frametitle{Regla de Decisión por Video}

  \begin{columns}[T]
    \column{0.31\textwidth}
      \metriccard{Chunk}{accentblue}{Etiqueta}
      \vspace{1mm}
      {\small Cada fragmento recibe categoría, probabilidad y texto activador.}

    \column{0.31\textwidth}
      \metriccard{Video}{darkred}{Riesgo}
      \vspace{1mm}
      {\small Se marca si varios chunks cruzan umbral o uno crítico es muy alto.}

    \column{0.31\textwidth}
      \metriccard{Salida}{accentteal}{Evidencia}
      \vspace{1mm}
      {\small Tabla con video, categoría, score y fragmentos revisables.}
  \end{columns}

  \vspace{5mm}
  \begin{block}{Principio}
    La decisión automática prioriza revisión humana; no sanciona contenido por sí sola.
  \end{block}
\end{frame}

% ==================== 9 ====================
\begin{frame}
  \frametitle{Cuadernos Iniciales}

  \begin{center}
  \begin{tikzpicture}[
    card/.style={draw=linegray,rounded corners=8pt,fill=white,minimum width=2.35cm,minimum height=1.35cm,align=center,font=\scriptsize},
    arr/.style={-Latex,thick,draw=muted}
  ]
    \node[card,fill=softblue] (n1) {\textbf{01}\\Scraping\\\texttt{raw.jsonl}};
    \node[card,right=0.25cm of n1,fill=softgreen] (n2) {\textbf{02}\\Limpieza\\chunks};
    \node[card,right=0.25cm of n2,fill=softgold] (n3) {\textbf{03}\\Etiquetado\\HTML};
    \node[card,right=0.25cm of n3,fill=softred] (n4) {\textbf{04}\\Entrenamiento\\modelo};
    \node[card,right=0.25cm of n4,fill=softgray] (n5) {\textbf{05}\\Producción\\demo};
    \draw[arr] (n1) -- (n2);
    \draw[arr] (n2) -- (n3);
    \draw[arr] (n3) -- (n4);
    \draw[arr] (n4) -- (n5);
  \end{tikzpicture}
  \end{center}

  \vspace{4mm}
  \begin{columns}[T]
    \column{0.48\textwidth}
      \begin{block}{Entrada y curaduría}
        \texttt{yt-dlp} y Selenium generan metadatos, subtítulos y chunks trazables. No se usan APIs ni etiquetas automáticas externas.
      \end{block}

    \column{0.48\textwidth}
      \begin{block}{Salida del artefacto}
        El flujo produce dataset etiquetado, baseline local, métricas y frontend de inferencia para revisión humana.
      \end{block}
  \end{columns}
\end{frame}

% ==================== 10 ====================
\begin{frame}
  \frametitle{Modelo y Entrenamiento}

  \begin{columns}[T]
    \column{0.48\textwidth}
      \begin{block}{Baseline obligatorio}
        \begin{itemize}
          \item TF-IDF con unigramas y bigramas.
          \item Linear SVM o Logistic Regression.
          \item One-vs-rest si se conserva multi-etiqueta.
          \item Explicable mediante términos con mayor peso.
        \end{itemize}
      \end{block}

    \column{0.48\textwidth}
      \begin{block}{Extensión avanzada}
        \begin{itemize}
          \item BETO o RoBERTa en español.
          \item Fine-tuning solo si el dataset etiquetado es suficiente.
          \item Comparación contra baseline, no reemplazo directo.
        \end{itemize}
      \end{block}
  \end{columns}

  \vfill
  {\small\color{muted}\textbf{Criterio de defensa:} primero funciona el baseline; el Transformer solo se justifica si mejora F1 macro, recall de riesgo y análisis de errores.}
\end{frame}

% ==================== 11 ====================
\begin{frame}
  \frametitle{Evaluación}

  \begin{center}
  \begin{tikzpicture}[
    eval/.style={draw=linegray,rounded corners,fill=white,minimum width=2.65cm,minimum height=0.72cm,align=center,font=\scriptsize},
    arr/.style={-Latex,thick,draw=muted}
  ]
    \node[eval,fill=softblue] (a) {Acuerdo humano\\Cohen's kappa};
    \node[eval,right=0.45cm of a] (b) {Split\\70 / 15 / 15};
    \node[eval,right=0.45cm of b] (c) {F1 macro\\recall por riesgo};
    \node[eval,right=0.45cm of c,fill=softgreen] (d) {Errores\\y umbrales};
    \draw[arr] (a) -- (b);
    \draw[arr] (b) -- (c);
    \draw[arr] (c) -- (d);
  \end{tikzpicture}
  \end{center}

  \vspace{3mm}
  \begin{columns}[T]
    \column{0.48\textwidth}
      \begin{block}{Métricas}
        \begin{itemize}
          \item F1 macro para clases desbalanceadas.
          \item Precision y recall por categoría.
          \item Matriz de confusión.
          \item Precisión de top-fragmentos revisados.
        \end{itemize}
      \end{block}

    \column{0.48\textwidth}
      \begin{block}{Calidad humana}
        \begin{itemize}
          \item Dos anotadores para una muestra.
          \item Revisión de desacuerdos.
          \item Guía de etiquetas con ejemplos.
          \item Registro de casos ambiguos.
        \end{itemize}
      \end{block}
  \end{columns}
\end{frame}

% ==================== 12 ====================
\begin{frame}
  \frametitle{Estructura del Paper / Informe}

  \begin{center}
  \begin{tikzpicture}[
    sec/.style={draw=linegray,rounded corners,fill=white,minimum width=2.7cm,minimum height=0.72cm,align=center,font=\scriptsize},
    arr/.style={-Latex,thick,draw=darkred}
  ]
    \node[sec,fill=softblue] (s1) {1. Contexto\\y problema};
    \node[sec,right=0.35cm of s1] (s2) {2. Brecha\\y objetivos};
    \node[sec,right=0.35cm of s2] (s3) {3. Método\\y corpus};
    \node[sec,below=0.65cm of s3,fill=softgreen] (s4) {4. Artefacto\\y notebooks};
    \node[sec,left=0.35cm of s4] (s5) {5. Evaluación\\y resultados};
    \node[sec,left=0.35cm of s5,fill=softred] (s6) {6. Discusión\\y límites};
    \draw[arr] (s1) -- (s2);
    \draw[arr] (s2) -- (s3);
    \draw[arr] (s3) -- (s4);
    \draw[arr] (s4) -- (s5);
    \draw[arr] (s5) -- (s6);
  \end{tikzpicture}
  \end{center}

  \begin{block}{Hilo conductor}
    Lenguaje local en videos peruanos $\rightarrow$ riesgo textual difícil de revisar manualmente $\rightarrow$ corpus anotado $\rightarrow$ clasificador local $\rightarrow$ evidencia para moderación humana.
  \end{block}
\end{frame}

% ==================== 13 ====================
\begin{frame}
  \frametitle{Alcances y Límites}

  \begin{columns}[T]
    \column{0.48\textwidth}
      \begin{block}{Alcance}
        \begin{itemize}
          \item Moderación textual, no visual.
          \item Videos públicos de política y farándula.
          \item Chunks de transcripción o subtítulos.
          \item Clasificación local con evidencia.
        \end{itemize}
      \end{block}

    \column{0.48\textwidth}
      \begin{block}{Límites}
        \begin{itemize}
          \item No afirma verdad o falsedad absoluta.
          \item No sustituye evaluación legal ni editorial.
          \item Puede fallar con ironía, sarcasmo y citas.
          \item Requiere balance y revisión de sesgos.
        \end{itemize}
      \end{block}
  \end{columns}

  \vspace{-1mm}
  \begin{block}{Principio ético}
    La salida del sistema es alerta y priorización, no sanción automática.
  \end{block}
\end{frame}

% ==================== 14 ====================
\begin{frame}
  \frametitle{Plan de Trabajo}

  \begin{columns}[T]
    \column{0.24\textwidth}
      \metriccard{Semana 1}{accentblue}{Diseño}
      \vspace{1mm}
      {\small Canales, categorías MVP, guía de anotación y criterios de descarte.}

    \column{0.24\textwidth}
      \metriccard{Semana 2}{accentteal}{Corpus}
      \vspace{1mm}
      {\small Scraping, subtítulos, transcripción local y creación de chunks.}

    \column{0.24\textwidth}
      \metriccard{Semana 3}{accentgold}{Etiquetado}
      \vspace{1mm}
      {\small Anotación humana, revisión de desacuerdos y dataset final.}

    \column{0.24\textwidth}
      \metriccard{Semana 4}{darkred}{Demo}
      \vspace{1mm}
      {\small Baseline, métricas, análisis de errores y frontend de producción.}
  \end{columns}

  \vspace{6mm}
  \begin{block}{Mejora propuesta}
    Mantener un registro de decisiones: por qué se incluye cada canal, por qué se descarta cada video y qué regla resolvió cada caso ambiguo.
  \end{block}
\end{frame}

% ==================== 15 ====================
\begin{frame}
  \frametitle{Referencias y Fuentes Base}

  \scriptsize
  \begin{thebibliography}{99}
    \bibitem[Radford et al., 2022]{radford2022whisper}
    Radford, A. et al. (2022).
    \newblock Robust Speech Recognition via Large-Scale Weak Supervision.
    \newblock \textit{arXiv:2212.04356}.

    \bibitem[Devlin et al., 2018]{devlin2018bert}
    Devlin, J. et al. (2018).
    \newblock BERT: Pre-training of Deep Bidirectional Transformers for Language Understanding.
    \newblock \textit{arXiv:1810.04805}.

    \bibitem[Cañete et al., 2020]{canete2020beto}
    Cañete, J. et al. (2020).
    \newblock Spanish Pre-Trained BERT Model and Evaluation Data.
    \newblock \textit{PML4DC at ICLR}.

    \bibitem[Ribeiro et al., 2020]{ribeiro2020checklist}
    Ribeiro, M. T. et al. (2020).
    \newblock Beyond Accuracy: Behavioral Testing of NLP Models with CheckList.
    \newblock \textit{ACL}, 4902--4912.

    \bibitem[Fuentes web, 2026]{fuentes2026}
    Canales públicos de YouTube revisables: RPP Noticias, TVPerú Noticias, El Comercio, La República, La Encerrona, Epicentro TV, Magaly TV La Firme, Amor y Fuego, América Hoy.

    \bibitem[Zavala \& Back, 2017]{zavala2017racismo}
    Zavala, V. \& Back, M. (Eds.) (2017).
    \newblock \textit{Racismo y lenguaje}. Fondo Editorial PUCP. 409 pp.

    \bibitem[Zavala \& Zariquiey, 2007]{zavala2007discurso}
    Zavala, V. \& Zariquiey, R. (2007).
    \newblock ``Yo te segrego a ti...'': discurso racista en el Perú contemporáneo.
    \newblock En Wodak \& van Dijk (Eds.), \textit{Racismo y discurso en América Latina}, pp.~243--285. Gedisa.

    \bibitem[Almeida \& Zavala, 2022]{almeida2022motoso}
    Almeida, C. \& Zavala, V. (2022).
    \newblock ``Motoso y terruco'': ideologías lingüísticas y racialización en la política peruana.
    \newblock \textit{Lexis}, 46(2). doi:10.18800/lexis.202202.004.

    \bibitem[Brañez Medina, 2012]{branez2012amixer}
    Brañez Medina, R. F. (2012).
    \newblock Identidades ``amixer'' y ``no-amixer'': racismo cultural en el espacio virtual.
    \newblock \textit{Discurso \& Sociedad}, 6(1), 1--35.

    \bibitem[Callirgos, 1993]{callirgos1993racismo}
    Callirgos, J. C. (1993).
    \newblock \textit{El racismo: la cuestión del otro (y de uno)}. DESCO, Lima.

    \bibitem[Portocarrero, 2009]{portocarrero2009racismo}
    Portocarrero, G. (2009).
    \newblock \textit{Racismo y mestizaje y otros ensayos}. Fondo Editorial del Congreso del Perú.

    \bibitem[Vich, 2018]{vich2018dinamicas}
    Vich, V. (2018).
    \newblock Dinámicas de racismo en el Perú: la perspectiva de Portocarrero.
    \newblock \textit{Debates en Sociología}, (47), 127--145.

    \bibitem[Monge-Olivarría \& Guerra-Corrales, 2023]{monge2023violencia}
    Monge-Olivarría, C. \& Guerra-Corrales, J. (2023).
    \newblock Violencia de género en Twitter: feminización como insulto.
    \newblock \textit{Lengua, Literatura y Arte}.

    \bibitem[Thakur / CDT, 2025]{thakur2025quechua}
    Thakur, D. (2025).
    \newblock \textit{Moderación de contenido en quechua en redes sociales}. CDT.
  \end{thebibliography}
\end{frame}

\end{document}
